%%%%%%%%%%%%%%%%%%%%%%%%%%%%%%%%%%%%%%%%%%%%%%%%%%%%%%%%%%%%%%%%%%%%%%
% Lớp 4 | Chủ đề 1. Ôn tập và bổ sung | Biểu thức chứa chữ
% Dạng 4D1?4-1: Tính giá trị biểu thức chứa chữ
%%%%%%%%%%%%%%%%%%%%%%%%%%%%%%%%%%%%%%%%%%%%%%%%%%%%%%%%%%%%%%%%%%%%%%


% nguồn: 2026-11-18_lop4_chu-de-1-tu-luan (nhập 18/11/2026)
\begin{bt}%[4D1N4-1]%[ma:d3341dfe]
% nguon: Toán tư duy lớp 4 – Bài 4 SGK Toán 4 KNTT
	Tính giá trị của biểu thức $125 + a \times 4$ với $a = 30$.
	\loigiai{
		Thay $a = 30$: $125 + 30 \times 4 = 125 + 120 = 245$.\\
		Nhớ: làm phép nhân trước, phép cộng sau.
	}
\end{bt}

% nguồn: 2026-11-18_lop4_chu-de-1-tu-luan (nhập 18/11/2026)
\begin{bt}%[4D1H4-1]%[ma:50724325]
% nguon: Toán tư duy lớp 4 – Bài 4 SGK Toán 4 KNTT
	Tính giá trị của biểu thức $(a + b) \times c$ với $a = 18$, $b = 32$, $c = 7$.
	\loigiai{
		Thay số: $(18 + 32) \times 7 = 50 \times 7 = 350$.\\
		Nhớ: làm phép tính trong ngoặc trước.
	}
\end{bt}

% nguồn: 2026-11-18_lop4_chu-de-1-tu-luan (nhập 18/11/2026)
\begin{bt}%[4D1V4-1]%[ma:8cf6fc5d]
% nguon: Toán tư duy lớp 4 – Bài 4 SGK Toán 4 KNTT
	Cho biểu thức $A = 500 - m \times 8$ ($m$ là số tự nhiên).\\
	a) Tìm $m$ để $A = 300$.\\
	b) Với $m$ bằng bao nhiêu thì $A$ có giá trị nhỏ nhất? Giá trị nhỏ nhất đó là bao nhiêu?
	\loigiai{
		a) $500 - m \times 8 = 300$ nên $m \times 8 = 500 - 300 = 200$, $m = 200 : 8 = 25$.\\
		b) $A$ nhỏ nhất khi $m \times 8$ lớn nhất nhưng không vượt quá $500$.\\
		$500 : 8 = 62$ (dư $4$), nên $m \times 8$ lớn nhất là $62 \times 8 = 496$.\\
		Vậy $m = 62$ và giá trị nhỏ nhất của $A$ là $500 - 496 = 4$.
	}
\end{bt}

% nguồn: 2026-11-18_lop4_chu-de-1-tu-luan-bo-sung (nhập 18/11/2026)
\begin{bt}%[4D1H4-1]%[ma:3cdf8ce4]
% nguon: Toán tư duy lớp 4 – Bài 4 SGK Toán 4 KNTT
	Tính giá trị của biểu thức $a \times b - c : 4$ với $a = 125$, $b = 8$, $c = 400$.
	\loigiai{
		Thay số: $125 \times 8 - 400 : 4$.\\
		Làm nhân, chia trước: $125 \times 8 = 1\,000$; $400 : 4 = 100$.\\
		Rồi trừ: $1\,000 - 100 = 900$.
	}
\end{bt}
